\documentclass[11pt]{article}
\usepackage[margin=2.6cm]{geometry}
\usepackage{mathpazo}
\usepackage{microtype}
\usepackage{booktabs}
\usepackage{graphicx}
\usepackage{float}
\usepackage{xcolor}
\usepackage{titlesec}
\usepackage{caption}
\usepackage{listings}
\usepackage[hidelinks]{hyperref}
\emergencystretch=3em

\definecolor{accent}{HTML}{1F4E79}
\definecolor{good}{HTML}{1A7F37}
\definecolor{bad}{HTML}{B42318}
\definecolor{warn}{HTML}{B7791F}
\titleformat{\section}{\Large\bfseries\color{accent}}{\thesection}{0.8em}{}
\titleformat{\subsection}{\large\bfseries\color{accent!80!black}}{\thesubsection}{0.8em}{}
\captionsetup{font=small,labelfont={bf,color=accent}}
\newcommand{\armA}{\textcolor{bad}{arm A}}
\newcommand{\armB}{\textcolor{accent}{arm B}}

\lstdefinestyle{json}{
  basicstyle=\ttfamily\footnotesize,
  backgroundcolor=\color{gray!8},
  frame=single, rulecolor=\color{gray!40}, framerule=0.4pt,
  breaklines=true, columns=fullflexible, keepspaces=true,
  showstringspaces=false, aboveskip=8pt, belowskip=8pt}
\lstset{style=json}

\newsavebox{\eliBox}
\newenvironment{eli5}
  {\par\medskip\noindent\begin{lrbox}{\eliBox}\begin{minipage}{\dimexpr\linewidth-2\fboxsep-2\fboxrule}\small}
  {\end{minipage}\end{lrbox}\noindent\fcolorbox{accent!40}{accent!4}{\usebox{\eliBox}}\par\medskip}

\newenvironment{interimbox}
  {\par\medskip\noindent\begin{lrbox}{\eliBox}\begin{minipage}{\dimexpr\linewidth-2\fboxsep-2\fboxrule}\small}
  {\end{minipage}\end{lrbox}\noindent\fcolorbox{warn!60}{warn!8}{\usebox{\eliBox}}\par\medskip}

\title{\color{accent}\bfseries Tool-Aware DAG-Diff Mutation for Retrieval Chains\\
\large Does the structured diff still earn its keep when the mutator is a
\emph{top-tier reasoning model}?\\
\large \textcolor{warn}{\bfseries INTERIM} --- live partial A/B, single-parent,
\texttt{Qwen3-235B-A22B-Thinking} mutator}
\author{GigaEvo}
\date{2026-07-04}

\begin{document}
\maketitle

\begin{interimbox}
\textbf{\textcolor{warn}{INTERIM SNAPSHOT --- both arms are still running.}} Numbers
are a single consistent snapshot of two live runs (target 250 mutants/arm; at snapshot
\armA{} has 185 valid + 6 in-flight children, \armB{} 179 valid + 5 in-flight). Every
figure and table in this report was regenerated from the \emph{same} storage/log pull, so
they agree with one another, but all values will move as the runs complete. No held-out
test evaluation has been run on these live programs yet, so the gemini report's held-out
coverage table and lineage walk-through are intentionally \emph{absent} here, not stubbed.
\end{interimbox}

\begin{abstract}
\noindent
This is the companion run to the shipped \texttt{gemini-3.5-flash} A/B: the identical
tool-aware slot-diff experiment on \textbf{HoVer multi-hop retrieval}
(\texttt{problems/chains/hover/full7}), but with the mutator swapped for a
\emph{top-tier reasoning model}, \texttt{Qwen3-235B-A22B-Thinking}, served over the LiteLLM
proxy. The gemini report left one question sharp: it argued the malformed-JSON failure of
free-form rewriting is \emph{structural} (weak \texttt{llama} malformed $\sim$40\%, gemini
38\% --- ``not model IQ''). The strongest reasoning mutator we have puts a number on that
claim, and the number is a partial rebuttal: over completed evaluations \armA{} (free
rewrite) malforms \textbf{24 of 209 (11.5\%)} of its wire-JSON children --- \emph{far} below
gemini's 38\%, so model quality \emph{does} cut the failure rate roughly threefold --- while
\armB{} (tool diff) malforms \textbf{0 of 179}. So the honest reading is: a stronger mutator
shrinks the malformed class but never closes it; only the schema guarantees zero. The
diff's \emph{value} is therefore largest against weak/mid mutators and narrows --- but does
not vanish --- as the mutator strengthens. But validity is not the diff's only genuine win,
and it would be unfair to the diff to file the rest under ``parity.'' \textbf{Where \armB{}
is outright better:} it wastes \textbf{\$0.00} of its budget on invalid children (vs \armA's
\$0.33); it ships \textbf{18\% fewer input tokens} (slot deltas, not a re-quoted document);
every child is a \textbf{typed, auditable diff of a known parent} rather than an opaque
re-emission; and --- crucially --- it reaches the \emph{same fitness ceiling} while exploring
a \textbf{much wider structural range} (2--7 step chains vs \armA's near-total collapse onto
7 steps). Best fitness is a genuine \emph{tie} (0.821 \armA{} vs 0.819 \armB{}, a 0.002 gap
inside the \texttt{Qwen3-8B} executor's $\sim$0.01--0.02 re-score noise), so the diff gives up
nothing at the top. \armA's lead is confined to two places, and both come with an asterisk:
the \emph{mean} (0.763 vs 0.714) is not a quality edge but the price of B's wider exploration
into low-scoring short chains --- i.e.\ B is doing \emph{more} search for the same ceiling ---
and \emph{cost-to-winner} (\$0.68 vs \$1.02), where the diff's dramatic gemini-era win (\$6.46
vs \$15.93) shrinks because \armA{} now wastes only 11\% (vs 38\%) so total spend converges
(\$2.88 vs \$2.98). One methodological bonus over the gemini run: the LiteLLM proxy does not
itemise a reasoning channel for \emph{either} arm (both log 0 reasoning tokens), so there is
no one-sided instrumentation confound here --- total billed output is a clean cross-arm read,
and it too is near-parity (\armB{} output $+4.9\%$). Finally we document the diff's \emph{one}
residual grammar failure
class --- \texttt{diff\_schema\_error} (slot-contiguity, a cross-field validator that
constrained decoding cannot cover) --- of which this run produced \textbf{zero}, and the new
deterministic slot-compaction repair that now closes even that class.
\end{abstract}

\section{Motivation}
CARL chain genomes are wire-JSON documents. The stock mutation path
(\texttt{mutation=llm\_rewrite}) asks the LLM to re-emit the \emph{whole} document as free
text; a malformed child is discovered only \emph{after} it has been persisted and has
consumed a full evaluation. The structured slot-diff (\texttt{mutation=carl\_with\_%
retrieval\_tools}) instead ships a per-mutation JSON schema so a malformed child cannot be
sampled at all. Two prior A/Bs established this on the LLM-only summarizer task; the shipped
\texttt{gemini-3.5-flash} run carried it to tool-bearing HoVer chains and found the validity
gap holds even against a strong mutator (61.6\% vs 98.8\% valid). This run isolates the last
free variable those left open:
\begin{itemize}
  \item \textbf{Does the guarantee still pay when the mutator is \emph{the strongest reasoning
  model available}?} The gemini report's central claim was that malformed JSON is structural
  to full-document re-emission, not a weak-model artifact --- evidenced by weak \texttt{llama}
  ($\sim$40\%) and \texttt{gemini-3.5-flash} (38\%) landing at nearly the same rate. If that
  is right, \texttt{Qwen3-235B-A22B-Thinking} --- a far more capable model --- should still
  malform $\sim$38\%. If instead the rate drops sharply, the failure \emph{is} partly a model
  effect, and the diff's value is a function of \emph{which} mutator you pair it with. That is
  the question this snapshot answers.
\end{itemize}

\section{Setup}
\label{sec:setup}
Two arms, run from scratch, identical except the mutation operator:
\begin{itemize}
  \item \textbf{Task:} \texttt{problems/chains/hover/full7}, \texttt{pipeline=program\_is\_json}.
  Full-chain mode (no frozen steps, no topology matching): the child may freely change step
  count/types/dependencies within \texttt{FULL\_CHAIN\_CONFIG} (\texttt{max\_steps}=7). Fitness
  is soft retrieval coverage --- for each of 300 HoVer claims, the fraction of gold
  supporting-fact article titles found across all tool-step outputs (higher better).
  \item \textbf{Mutator:} \texttt{Qwen3-235B-A22B-Thinking} via the LiteLLM proxy at
  \texttt{10.232.89.98:4000} --- \emph{identical config on both arms}. This is the same mutator
  family as the user's HoVer feedback-condition sweep, so the mutator axis here is comparable
  to that prior work.
  \item \textbf{Executor:} \texttt{Qwen/Qwen3-8B} via the same proxy, step-batched over the
  300-sample set --- \emph{identical to the gemini run}, so the fitness scale is shared.
  \item \textbf{Config:} \texttt{memory=none}, \texttt{num\_parents=1}, one 7-step retrieval
  seed, disk storage, target \texttt{max\_mutants=250}. \armA{} = \texttt{mutation=llm\_rewrite}
  (stock full wire-JSON rewrite); \armB{} = \texttt{mutation=carl\_with\_retrieval\_tools}
  (\texttt{StructuredDiffMutationOperator} + \texttt{AllowedToolChainChanges}).
\end{itemize}

\noindent\textbf{Single parent.} With \texttt{num\_parents=1} the diff grammar collapses to
its single-parent form; crossover is unavailable this run (a deliberate simplification, matching
the gemini run for comparability).

\section{The tool-aware diff language}
\label{sec:language}
The child chain is a shelf of numbered slots the LLM fills left-to-right; a slot may hold a
\texttt{tool} step (a \texttt{retrieve}/\texttt{retrieve\_deep} with a named query source) as well
as an LLM step. This is unchanged from the gemini run --- the \emph{grammar} is fixed; only the
mutator model differs --- so we keep the description short.

\begin{eli5}
Picture the child chain as a row of numbered boxes. Into each box you put either a
\emph{photocopy} of a step from the parent (named by its label \texttt{a1}, \texttt{a2}, \dots,
optionally with small corrections) or a \emph{brand-new} step; ``go-look-it-up'' \emph{tool}
steps are photocopied the same way. Each box lists which \emph{earlier} boxes it reads from,
and the form for box $k$ only has checkboxes for boxes $1..k{-}1$ --- so pointing forward or at
yourself is not against the rules, there is simply nowhere on the form to write it. Leftover
boxes stay empty (no gaps), and the last filled box produces the graded answer.
\end{eli5}
\noindent
The form is a JSON schema compiled fresh per mutation and shipped as the provider's
\texttt{json\_schema} structured-output grammar, so a violating token can never be sampled; a
pydantic model re-checks as a backstop and the applier round-trips through \texttt{ReasoningChain}
as a final assertion. Everything an unconstrained rewriter routinely gets wrong ---
self/forward dependencies, cycles, keeping a nonexistent step, invented fields, \emph{malformed
JSON} --- is unrepresentable by construction; only step-contiguity is left to a validator
(\S\ref{sec:taxonomy}). A rejection costs one LLM call: nothing persisted, nothing evaluated.

\section{Results}
\label{sec:results}

\subsection{Validity funnel --- the headline, and the model-quality twist}
\begin{table}[H]\centering
\begin{tabular}{lcc}
\toprule
live snapshot (target 250/arm) & \armA{} (rewrite) & \armB{} (tool diff) \\
\midrule
mutation LLM calls (incl.\ retries) & 215 & 191 \\
pre-persist generation failures & 7 (no code) & 9 (LLM-call) \\
children persisted & 215 & 184 \\
\quad still executing at snapshot (excluded) & 6 & 5 \\
\midrule
\textbf{completed evaluations} & \textbf{209} & \textbf{179} \\
\quad valid & 185 & 179 \\
\quad invalid & 24 & \textbf{0} \\
\qquad of which malformed JSON & \textbf{\textcolor{warn}{24 (11.5\%)}} & \textbf{\textcolor{good}{0}} \\
\textbf{valid yield among completed} & \textbf{\textcolor{warn}{88.5\%}} & \textbf{\textcolor{good}{100.0\%}} \\
\bottomrule
\end{tabular}
\caption{The diff still removes the malformed class entirely (\textbf{0} malformed-JSON children),
but the free-form path is much healthier here than under gemini: \textbf{11.5\% malformed} vs
gemini's 38\%. \armA's remaining pre-persist failures are 7 \texttt{no\_code\_extracted} (the model
returned prose with no extractable JSON); \armB's are 9 \texttt{llm\_call\_error} --- proxy/timeout
faults, \emph{not} grammar violations (see \S\ref{sec:taxonomy}). Both arms exclude in-flight
children (the ``unknown'' invalids in \texttt{ab\_stats.json} are exactly these not-yet-scored
programs: \armA's 6, \armB's 5).}
\end{table}
\begin{figure}[H]\centering
\includegraphics[width=\textwidth]{ab_overview.png}
\caption{Left: per-child soft-coverage fitness (dots) and best-so-far (line); invalid children are
crosses on the floor line --- present for \armA{} (\textcolor{bad}{red}), \emph{entirely absent} for
\armB{}. The two best-so-far lines are on top of each other (\armA{} ends marginally higher, 0.821 vs
0.819). Note \armB's fatter low tail (dots down to $\sim$0.63): those are its short-chain explorations,
which drag its \emph{mean} down (\S\ref{sec:struct}). Right: the validity funnel --- both persist
$\sim$190--215, \armA{} keeps 185 valid, \armB{} keeps 179.}
\end{figure}

\noindent\textbf{Why 11.5\%, not 38\%.} This is the run's most consequential number, and it
qualifies the gemini report. That report argued the malformed rate is structural to
full-document re-emission and roughly model-independent (weak \texttt{llama} $\sim$40\%, gemini
38\%). \texttt{Qwen3-235B-A22B-Thinking} --- a materially stronger reasoning model --- malforms
\emph{roughly a third} as often (11.5\%). So the failure is \emph{partly} a model effect: a
top-tier mutator is genuinely better at re-emitting a whole tool-bearing chain as valid JSON.
What survives of the gemini claim is the weaker, correct version: even the strongest mutator
does \emph{not} reach zero, and the gap it leaves ($1$-in-$9$ children wasted on a malformed
re-emission, each having burned an evaluation) is exactly the class the schema removes outright.
The diff's validity value is thus a \emph{function of the mutator}: decisive against weak/mid
models, a modest-but-real safety margin against the strongest.

\subsection{Fitness --- the ceiling is a tie; \armA's mean lead is B's exploration tax}
\begin{table}[H]\centering
\begin{tabular}{lccc}
\toprule
completed children & \armA{} (rewrite) & \armB{} (tool diff) & read \\
\midrule
best child & 0.821 & 0.819 & \textbf{tie ($\Delta$0.002 $<$ noise)} \\
child mean & \textbf{0.763} & 0.714 & A leads (B explores more) \\
child median & \textbf{0.773} & 0.758 & A leads (same reason) \\
seed (executor re-score) & 0.750 & 0.743 & tie \\
\bottomrule
\end{tabular}
\caption{The number that matters --- the \emph{best} program each arm can reach --- is a dead heat
(0.821 vs 0.819, a 0.002 gap well inside the \texttt{Qwen3-8B} executor's $\sim$0.01--0.02 re-score
noise). So the diff gives up \emph{nothing} at the ceiling. \armA's lead on the \emph{mean/median}
is not a quality edge: it is the visible cost of B doing more work --- the diff arm spends a slice of
its budget on short-chain structural experiments (\S\ref{sec:struct}) that score low, which pulls its
\emph{average} down without touching its \emph{ceiling}. Framed for the diff: B reaches the same top
fitness \emph{and} searches a wider structural neighbourhood to get there. Read honestly, neither arm
buys a higher optimum on this task; the diff's case is that it ties the optimum while winning on
validity, waste, input cost, legibility, and exploration breadth.}
\end{table}

\subsection{Structural exploration --- the diff arm ranges wider}
\label{sec:struct}
\begin{table}[H]\centering
\begin{tabular}{lccccccc}
\toprule
valid children by chain length & 2 & 3 & 4 & 5 & 6 & 7 & \% at 7 \\
\midrule
\armA{} (rewrite)   & --- & 1 & 5 & 7 & 14 & 158 & \textbf{85\%} \\
\armB{} (tool diff) & 3 & 3 & 8 & 12 & 21 & 132 & \textbf{74\%} \\
\bottomrule
\end{tabular}
\caption{The reversal from the gemini run. There, the diff stayed \emph{tighter} to the 7-step seed
than free rewrite (84\% vs 77\% at 7 steps). Here it is the opposite: \armB{} ranges more widely
toward short chains (26\% off-7, including 2- and 3-step chains free rewrite never produced), while
\armA{} clings to 7 steps (85\%) --- a near-mode-collapse onto the seed's shape. The diff's cheap,
always-valid short-chain edits (delete-to-2-steps is one representable slot move) let the strong
mutator \emph{afford} structural experiments free rewrite never even attempts. On HoVer full7 those
short chains mostly underperform (hence \armB's lower mean), so the breadth does not lift the ceiling
\emph{here} --- but structural diversity at a tied ceiling is a genuine property of the diff, not a
defect: it is exactly the behaviour you want on a harder task, a larger budget, or a landscape where
the seed's 7-step shape is \emph{not} already near-optimal. Free rewrite's tidy 85\%-at-7 is the flip
side of the same coin: less waste-to-explore, but also less willingness to leave the seed's basin.}
\end{table}
\begin{figure}[H]\centering
\includegraphics[width=0.72\textwidth]{ab_nsteps.png}
\caption{Chain-length distribution of valid children. \armB{} (blue) spreads across 2--7 steps;
\armA{} (red) concentrates at 7.}
\end{figure}

\subsection{Token cost --- no reasoning confound this time}
\label{sec:tokens}
The gemini run had a one-sided instrumentation bug: only the diff agent logged a reasoning-token
channel, so its total-output comparison came with a caveat. \textbf{That confound is absent here}
--- the LiteLLM proxy reports \emph{no} reasoning channel for either arm (both log 0 reasoning
tokens), so \texttt{tokens\_out} is total generated output on both, measured identically, and the
cross-arm comparison is clean.
\begin{table}[H]\centering
{\small
\begin{tabular}{lccc}
\toprule
mutation side (\texttt{Qwen3-235B}), snapshot & \armA{} & \armB{} & \armB{} vs \armA{} \\
\midrule
tokens in & 1{,}260{,}347 & 1{,}030{,}499 & $-18.2\%$ \\
\textbf{tokens out (total)} & \textbf{1{,}802{,}075} & \textbf{1{,}890{,}614} & \textbf{$+4.9\%$} \\
\quad of which reasoning & \multicolumn{2}{c}{\textcolor{warn}{not itemised by proxy (n/a both arms)}} & --- \\
mean output / completed call & 8{,}664 & 10{,}388 & $+20\%$ \\
LLM calls (incl.\ retries) & 215 & 191 & $-11\%$ \\
\bottomrule
\end{tabular}}
\caption{Near-parity on total output ($+4.9\%$ for the diff), and the diff sends \emph{18\% fewer}
input tokens (it ships slot deltas against a compact rendered parent rather than re-quoting a whole
document). Per accepted call the diff generates $\sim$20\% more (its structured payload plus the
model's un-itemised chain-of-thought), but it makes 11\% fewer calls, so totals land close. Because
the proxy does not split reasoning from content, we cannot state a content-only ratio --- but unlike
the gemini run, nothing here is one-sided.}
\end{table}
\begin{figure}[H]\centering
\includegraphics[width=\textwidth]{ab_tokens.png}
\caption{Left: mutation output tokens per call (dots, 10-call rolling mean) --- the diff's per-call
output runs slightly higher. Right: cumulative tokens; \armB's output (solid) edges above \armA's
while its input (dashed) sits below.}
\end{figure}

\subsection{Cost/quality frontier --- the diff's cost edge does not reproduce}
\label{sec:pareto}
Priced at OpenRouter's cheapest \texttt{Qwen3-235B-Thinking} rate (\$0.15/M in, \$1.495/M out;
the \texttt{Qwen3-8B} executor is a separate shared cost, excluded), the comparison that was
lopsided under gemini is \emph{near-parity} here --- and where it tilts, it tilts toward free
rewrite. Cost is attributed by \texttt{created\_at}: each child, in wall-clock creation order,
consumes the next accepted mutation call.
\begin{table}[H]\centering
\begin{tabular}{lccc}
\toprule
mutation-model \$ (OpenRouter cheapest) & \armA{} & \armB{} & \armB{} vs \armA{} \\
\midrule
total run spend (snapshot) & \$2.88 & \$2.98 & $+3.5\%$ \\
\midrule
cost to the best program so far & \textbf{\$0.68} & \$1.02 & $+50\%$ \\
\quad winner reached at (chrono.\ rank) & 52/215 (24\%) & 60/184 (33\%) & later \\
\quad winner fitness (train soft) & \textbf{0.821} & 0.819 & tie \\
\midrule
cost per \emph{valid} program & \textbf{\$0.016} & \$0.017 & $+6\%$ \\
budget wasted on invalid children & \textcolor{warn}{\$0.33 (11\%)} & \textbf{\textcolor{good}{\$0.00 (0\%)}} & --- \\
\bottomrule
\end{tabular}
\caption{Under gemini, the diff bought a \emph{higher} winner for \textbf{\$6.46 vs \$15.93}
cost-to-winner ($-59\%$) because \armA{} wasted 38\% of its budget on malformed children. That win
was a \emph{waste-differential} artifact: with \armA{} now wasting only 11\% (a strong mutator
malforms less), the differential shrinks and the arms converge. Here the diff is marginally
\emph{dearer} on every dollar except the one it is built to win --- wasted budget, where it is \$0.00
vs \$0.33. The clean statement: the diff's cost advantage scales with how much the free-rewrite arm
\emph{wastes}, which scales with how weak the mutator is; against the strongest mutator, free rewrite
wastes little and the diff's economic case narrows to its zero-waste guarantee.}
\end{table}
\begin{figure}[H]\centering
\includegraphics[width=0.86\textwidth]{ab_pareto.png}
\caption{Single-panel cost-to-winner vs.\ best train fitness (the gemini report's second panel,
held-out test coverage, is omitted --- no held-out eval has run on these live programs). Both arms'
cost-efficiency step lines climb to $\sim$0.82; \armA's open circle sits \emph{left} of \armB's
diamond (cheaper cost-to-winner) at marginally higher fitness. Neither arm dominates --- this is the
parity the strong mutator produces.}
\end{figure}
\begin{figure}[H]\centering
\includegraphics[width=\textwidth]{ab_waste.png}
\caption{\textbf{Left:} where each arm's budget went. \armA{} wastes \$0.33 of \$2.88 (11\%) on
malformed/invalid children; \armB{} wastes \$0.00. \textbf{Right:} effective cost per \emph{valid}
program --- \$0.016 (\armA) vs \$0.017 (\armB), a dead heat, because at 11\% waste the free-rewrite
penalty is small. Contrast the gemini run, where the same chart showed \$0.083 vs \$0.119 (a 30\%
diff advantage) driven by \armA's 38\% waste.}
\end{figure}

\section{Failure taxonomy and the deterministic contiguity repair}
\label{sec:taxonomy}
Every mutation that does not yield a scored child is classified by
\texttt{analyze\_ab.py}'s taxonomy (matched against the \texttt{MutationError} message). The
snapshot's full breakdown:
\begin{table}[H]\centering
\begin{tabular}{llcc}
\toprule
failure kind & where caught & \armA{} & \armB{} \\
\midrule
\texttt{json\_parse\_error} (malformed wire-JSON) & \textcolor{bad}{post-persist, after eval} & 24 & \textbf{\textcolor{good}{0}} \\
\texttt{no\_code\_extracted} (no JSON in reply) & pre-persist (free)          & 7  & \textbf{\textcolor{good}{0}} \\
\texttt{llm\_call\_error} (proxy / timeout)      & pre-persist (free)          & 0  & 9 \\
\texttt{diff\_schema\_error} (slot-contiguity)   & pre-persist (free)          & \textcolor{gray}{n/a} & \textbf{\textcolor{good}{0}} \\
in-flight (not yet scored)                        & ---                         & 6  & 5 \\
\bottomrule
\end{tabular}
\caption{The asymmetry is the whole point. \armA's failures are \emph{grammar/extraction}
(24 malformed + 7 no-code = 31 children that either wasted an evaluation or a call on badly-formed
output). \armB's only failures are \textbf{9 \texttt{llm\_call\_error}} --- infrastructure faults
(the proxy dropped the call), \emph{not} anything the model got wrong about the format. \armB{}
produced \textbf{zero} malformed JSON, \textbf{zero} extraction failures, and --- notably ---
\textbf{zero \texttt{diff\_schema\_error}}, the one grammar failure the structured diff can still hit.}
\end{table}

\subsection{The one thing constrained decoding cannot enforce}
Structured output guarantees the reply conforms to the JSON \emph{Schema}: types, enums, required
fields, per-field bounds. It cannot enforce a cross-field Python \texttt{@model\_validator}. The diff
has exactly one such rule --- \emph{slot contiguity} (``fill \texttt{slot\_1..slot\_k}
consecutively; no gaps''). It cannot be expressed per-field (no single slot's enum captures ``and
every earlier slot is also filled''), so the decoder is free to sample a schema-valid payload with a
hole (e.g.\ \texttt{slot\_3} filled, \texttt{slot\_2} null). Pydantic then rejects it, the operator
wraps it as \texttt{diff\_schema\_error}, and the mutation is retried --- one wasted call, nothing
persisted. Under gemini this fired on $2/259$ diffs (0.8\%); under \texttt{Qwen3-235B-Thinking} it
fired \textbf{0} times this snapshot. It is the benign end of the failure spectrum (caught pre-persist,
zero eval cost), but it is the diff's \emph{only} self-inflicted grammar failure, so it is worth
removing.

\subsection{The repair: slide filled slots left, remap references}
A gap diff is simply \emph{the right steps at the wrong indices}. The fix is deterministic
compaction, applied as a fallback the instant validation fails:
\begin{enumerate}
  \item Collect the filled slots in order, ignoring the holes.
  \item Slide them left into \texttt{slot\_1..slot\_m} (stable order preserved), building the
  old$\to$new index map.
  \item Rewrite every slot reference through that map --- each slot's backward
  \texttt{dependencies}, and (in the tool-chain variant) a tool step's \texttt{query\_source}.
  \item Re-validate the compacted payload. If it still fails --- e.g.\ a surviving slot depended on
  the \emph{hole} itself, which cannot be cleanly remapped --- abort the repair and let the original
  error stand.
\end{enumerate}
The move is an \emph{order-isomorphism}: because slot order is preserved and every edge already
points backward, compaction maps a gapped emission to \emph{byte-for-byte the chain a contiguous
emission would have produced}. Re-validating before use means the repair can never itself emit an
invalid child; the conservative abort means a genuinely ambiguous gap (a reference into the hole)
is still rejected rather than guessed. This is the same principle that makes dependency-adherence
decoder-enforceable, pushed one step further to cover the cross-field invariant the decoder cannot.
It is now wired into both diff vocabularies --- \texttt{problems/chains/tool\_chain\_diff.py}
(this run's operator) and its LLM-only cousin \texttt{gigaevo/chains/dag\_changes.py} --- each with
gap-repair and dangling-reference-rejection unit tests. So the \texttt{diff\_schema\_error} column
above is now not merely \emph{empty by luck of the mutator}, but \emph{empty by construction}: even
a model that does emit a slot gap will have it silently repaired into the valid chain it meant.

\section{Verdict}
Against a top-tier reasoning mutator (\texttt{Qwen3-235B-A22B-Thinking}), the structured tool-diff
still delivers its one hard guarantee --- \textbf{0 malformed-JSON children vs free rewrite's 11.5\%}
--- but the strong-mutator regime reframes what that guarantee is \emph{worth}. The headline finding
is a partial correction to the gemini report: malformed rate is \emph{not} model-independent. A
strong reasoning model malforms roughly a third as often as \texttt{gemini-3.5-flash} (11.5\% vs
38\%), so the diff's dramatic gemini-era \emph{margins} --- $+0.02$ fitness, $-59\%$ cost-to-winner,
$-30\%$ cost-per-valid --- shrink because \armA's waste shrinks. But it would be wrong to conclude
``the diff no longer wins.'' \textbf{Where \armB{} is genuinely better even against the strongest
mutator:} (i)~\emph{validity} --- 0 malformed vs 24, categorical; (ii)~\emph{waste} --- \$0.00 vs
\$0.33 of budget on invalid children; (iii)~\emph{input cost} --- 18\% fewer input tokens; (iv)~%
\emph{legibility} --- every child a typed, auditable diff of a named parent; (v)~\emph{exploration
breadth} --- it ties the fitness ceiling (0.821 vs 0.819, within noise) while ranging across 2--7
step chains where free rewrite collapses onto 7. \armA's only genuine wins are the \emph{mean} (which
is really B paying an exploration tax for that breadth) and \emph{cost-to-winner} (\$0.68 vs \$1.02,
a margin that exists only because \armA{} now wastes little). Net: the two arms tie on the optimum;
the diff wins the five properties above; free rewrite wins mean and cost-to-winner. The synthesis
across all three mutator runs is not ``the diff stops mattering'' but \textbf{the diff's \emph{fitness/cost}
margin scales inversely with mutator strength, while its validity/legibility/exploration advantages
hold at every strength} --- decisive against weak/mid models, a still-real safety-and-diversity margin
against the strongest.
Methodologically this run is also \emph{cleaner} than the gemini one: no reasoning-instrumentation
confound (the proxy itemises no reasoning channel for either arm), so total output is a fair
cross-arm read. \textbf{Caveats:} \textcolor{warn}{this is a live interim snapshot} (185/179 valid of
target 250, 6/5 in-flight excluded), one replicate per arm with no same-config variance floor yet on
HoVer full7 (so the fitness tie and mean gap should be read as \emph{provisional}, the validity gap
as categorical), \texttt{memory=none}, and \texttt{num\_parents=1} (no crossover). Next: let both arms
finish to 250; a variance floor on this task; and --- with the deterministic contiguity repair now
landed in both diff vocabularies --- confirm the \texttt{diff\_schema\_error} class stays at zero by
construction rather than by the mutator's good luck.
\end{document}
